% Краткое сообщение для УМН --- ДВУХСТРАНИЧНЫЙ вариант без результата в R^10.
% Оформление --- официальный шаблон (mathnet.ru/rus/rm/authornotes,
% umn_shap.zip: 000.tex + umnbib.sty), перекодированный в UTF-8.
% Отличия от note.tex: в теореме четыре оценки (43, 132, 1029, 7203), все ---
% точными рациональными сертификатами; п. 4 (продуктовое правило, планарная
% граница, 45619) удалён; ссылок на неопубликованные работы нет (10.09.2026):
% все данные четырёх решёток выписаны, матрицы 1029 и 7203 --- блочно. См. README.md.
\documentclass{article}
\usepackage[T2A]{fontenc}
\usepackage[utf8]{inputenc}
\usepackage[russian]{babel}
\usepackage[tbtags]{amsmath}
\usepackage{amsfonts,amssymb,mathrsfs,amscd,comment}
\usepackage{umnbib} % оформление списка литературы

\overfullrule10pt

\voffset-30mm\hoffset-15mm\mag1200
\textheight 230mm\textwidth 140mm\normalbaselineskip=12.5pt %4.39mm
\setlength{\emergencystretch}{1.5em}

\newcommand{\R}{\mathbb{R}}
\newcommand{\Z}{\mathbb{Z}}
\newcommand{\diam}{\operatorname{diam}}
\newcommand{\dist}{\operatorname{dist}}
\newtheorem{theorem}{Теорема}

\begin{document}

% TODO: сверить УДК с редакцией
\title{\makebox[\textwidth][l]{\normalfont\normalsize УДК 514.174.5+519.174}\\[8pt]
\bf\large\MakeUppercase{%
Новые верхние оценки хроматических чисел евклидовых пространств
}}

% Порядок авторов --- по вкладу (решение от 22.08.2026): в заметку входят
% решётка индекса 43 и точные верификаторы 1029 и 7203 (Н.\,В.~Глушкова).
\author{Л.\,Л.~Иванов, Н.\,В.~Глушкова}
\date{}

\maketitle

\makeatletter
\renewcommand{\@makefnmark}{}
\makeatother

\textbf{1. Постановка и результат.}
Через $\chi(\R^n,[1,\ell])$, $\ell\ge1$, обозначим наименьшее число цветов
раскраски пространства $\R^n$, при которой никакие две одноцветные точки не
находятся на расстоянии из отрезка $[1,\ell]$; при $\ell=1$ это классическое
хроматическое число $\chi(\R^n)$ (проблема Нельсона--Хадвигера, см.~\cite{Soifer});
интервальная постановка изучалась в~\cite{Ivanov2011}. Очевидно,
$\chi(\R^n)\le\chi(\R^n,[1,\ell])$. Лучшие известные верхние оценки в
размерностях $4$, $5$, $7$, $9$ --- соответственно $49$, $140$, $1372$ и
$17253$ \cite{ABPR}; в \cite{ABPR} высказано ожидание, что $49$ и $140$
неулучшаемы решёточными раскрасками.

\begin{theorem}
$\chi(\R^4)\le43$, $\chi(\R^5)\le132$, $\chi(\R^7)\le1029$,
$\chi(\R^9)\le7203$. Точнее, при всех $\ell\le\ell_0$ выполнено
$\chi(\R^n,[1,\ell])\le k$ для $(n,k,\ell_0)=(4,43,\,1{,}00411)$,
$(5,132,\,101/100)$, $(7,1029,\,103/100)$, а для $(n,k)=(9,7203)$ --- при всех
$\ell<1{,}0166127$.
\end{theorem}

Все четыре оценки даны явными рациональными решётками с проверкой в точной
арифметике (п.~3); код и полные протоколы проверок открыты.\footnote{\texttt{https://github.com/ivaleo/Chromatic}}

\textbf{2. Решёточные раскраски.}
Пусть $\Lambda\subset\R^n$ --- решётка, $\Gamma\subset\Lambda$ --- подрешётка
индекса $k$, $V_0$ --- ячейка Вороного нуля (выпуклый центрально-симметричный
многогранник); ячейку $a+V_0$, $a\in\Lambda$, окрасим цветом класса $a+\Gamma$.
Положим
\[
D(v)=2\dist\bigl(v/2,\,V_0\bigr),\qquad
D(\Gamma)=\min_{v\in\Gamma\setminus\{0\}}D(v),\qquad
d(\Lambda,\Gamma)=\frac{D(\Gamma)}{\diam V_0}.
\]
Если $x\in V_0$, $y\in v+V_0$, то $x'=v-y\in V_0$ и
$x-y=2\bigl(\tfrac{x+x'}2-\tfrac v2\bigr)$, откуда по выпуклости
$|x-y|\ge D(v)$. Значит,
две одноцветные точки либо лежат в одной ячейке (расстояние $\le\diam V_0$),
либо удалены не менее чем на $D(\Gamma)$. Если $d=d(\Lambda,\Gamma)>1$, то при
масштабе $s$ с $s\diam V_0<1$, $sD(\Gamma)>\ell$ (такой $s$ есть в точности при
$\ell<d$) отрезок $[1,\ell]$ свободен от одноцветных расстояний, то есть
\begin{equation}\label{crit}
\chi(\R^n,[1,\ell])\le[\Lambda:\Gamma]\quad\text{при всех }\ell<d(\Lambda,\Gamma).
\end{equation}

\textbf{3. Четыре явные конструкции.}
Пусть $Q$ --- рациональная положительно определённая матрица, $\Lambda=\Z^n$
со скалярным произведением $\langle x,y\rangle=x^{\mathsf T}Qy$ и
$\Gamma\subset\Lambda$ --- целочисленная подрешётка. Тогда неравенство
$d(\Lambda,\Gamma)>\ell_0$ сводится к конечной проверке в рациональных числах:
фасеты $V_0$ задаются релевантными векторами (по теореме Вороного они
находятся среди кратчайших представителей классов $\Lambda/2\Lambda$),
вершины $V_0$ --- решения рациональных линейных систем,
$\diam V_0=2R$, где $R^2$ --- максимум $\langle x,x\rangle$ по вершинам;
нарушить $D(v)\ge\ell_0\diam V_0$ может лишь вектор с $|v|<2(1+\ell_0)R$
(ибо $D(v)\ge|v|-\diam V_0$), таких конечное число, и для каждого
$\dist(v/2,V_0)$ подтверждается условиями Каруша--Куна--Таккера в дробях.

{\bfseries\boldmath $n=4$, $k=43$.} Решётка эйзенштейнова: $\Lambda=\Z[\omega]^2$,
$\omega=e^{2\pi i/3}$, с эрмитовой формой
$\bigl(\begin{smallmatrix}a&c\\\bar c&b\end{smallmatrix}\bigr)$, вещественная
часть которой в $\Z$-базисе $(e_1,\omega e_1,e_2,\omega e_2)$ при
$u=\operatorname{Re}c$, $s=\tfrac{\sqrt3}2\operatorname{Im}c$ имеет матрицу
Грама
\[
Q=\left(\begin{smallmatrix}
a & -a/2 & u & -u/2-s\\
-a/2 & a & -u/2+s & u\\
u & -u/2+s & b & -b/2\\
-u/2-s & u & -b/2 & b
\end{smallmatrix}\right),\qquad
a=1,\quad b=\tfrac{39261}{25000},\quad u=\tfrac{37381}{100000},\quad s=\tfrac{3141}{12500};
\]
подрешётка $\Gamma$ порождена строками матрицы
$T=\Bigl(\begin{smallmatrix}1&0&0&41\\0&1&0&14\\0&0&1&37\\0&0&0&43\end{smallmatrix}\Bigr)$,
$\det T=43$ ($43=N(7+\omega)$, $\Gamma$ --- $\Z[\omega]$-подмодуль). Ячейка
имеет $30$ фасет и $120$ вершин,
\[
\diam^2V_0=\tfrac{5374343884337}{2019775849050},\quad
D(\Gamma)^2=\tfrac{7396075258066147}{2756867811550000},\quad
d=1{,}0041105\ldots>1{,}00411 .
\]

{\bfseries\boldmath $n=5$, $k=132$.} Решётка общего положения:
\[
Q=\tfrac1{100000}\left(\begin{smallmatrix}
79327&86284&69597&36140&26501\\
86284&254538&216808&101257&102135\\
69597&216808&296568&163667&140527\\
36140&101257&163667&181733&71653\\
26501&102135&140527&71653&148076
\end{smallmatrix}\right),
\qquad
\Gamma=\ker\varphi,
\]
$\varphi(x)=88x_1+89x_2+127x_3+57x_4+3x_5 \bmod 132$ (гомоморфизм сюръективен,
$[\Lambda:\Gamma]=132$). Ячейка имеет $62$ фасеты и $720$ вершин,
{\small
\[
R^2=\tfrac{3093570095915736710467707313641}{3999974238630281363046165200000},\qquad
D(\Gamma)^2=\tfrac{1634114770527727}{516898614620000},\qquad
d=1{,}0108977\ldots>\tfrac{101}{100}.
\]}

{\bfseries\boldmath $n=7$, $k=1029$.} Ламинирование эйзенштейновой раскраски
$(3+\omega)E_6^*$ индекса $7^3=343$ тремя слоями: в подходящем $\Z$-базисе
(первые шесть векторов порождают базу $E_6^*$ с матрицей Грама $\tfrac34M$,
$\lambda_1^2=3$; седьмой --- слоевой)
$Q=\Bigl(\begin{smallmatrix}\frac34M&g\\ g^{\mathsf T}&17897/10^4\end{smallmatrix}\Bigr)$,
где
\[
M=\left(\begin{smallmatrix}
4&2&0&-3&2&1\\ 2&4&3&0&1&2\\ 0&3&6&3&0&3\\
-3&0&3&6&-3&0\\ 2&1&0&-3&4&2\\ 1&2&3&0&2&4
\end{smallmatrix}\right),\qquad
10^4g^{\mathsf T}=(-8632,\,-4197,\,1601,\,6617,\,-5275,\,4434)
\]
и $\Gamma=C\Z^7$, где $C$ --- верхнетреугольная матрица с диагональю
$(7,1,7,1,7,1,3)$ и единственными ненулевыми внедиагональными элементами
$C_{12}=C_{34}=C_{56}=-5$; $\det C=1029$. Ячейка имеет $254$ фасеты и
$30\,368$ вершин,
\[
\diam^2V_0=\tfrac{96858928789031597}{14761819462500000},\qquad
D(\Gamma)^2=7,\qquad
d=1{,}0328782\ldots>\tfrac{103}{100}.
\]

{\bfseries\boldmath $n=9$, $k=7203$.} То же ламинирование над $(3+\omega)E_8$ индекса
$7^4=2401$ с тремя слоями: в подходящем $\Z$-базисе (первые восемь векторов
порождают $E_8$, девятый --- слоевой)
$Q=\Bigl(\begin{smallmatrix}\frac32A&g\\ g^{\mathsf T}&17982/10^4\end{smallmatrix}\Bigr)$,
$\Gamma=C\Z^9$, где
\[
A=\left(\begin{smallmatrix}
2&1&1&-1&-1&1&1&1\\ 1&2&1&0&-1&1&1&1\\ 1&1&2&0&0&0&0&0\\ -1&0&0&2&0&0&0&-1\\
-1&-1&0&0&2&-1&-1&-1\\ 1&1&0&0&-1&2&1&1\\ 1&1&0&0&-1&1&2&1\\ 1&1&0&-1&-1&1&1&2
\end{smallmatrix}\right),\quad
C=\left(\begin{smallmatrix}
-5&-2&-9&-5&-3&-6&-10&-3&0\\ 4&3&-4&-3&0&0&-3&0&0\\ 0&0&9&5&3&3&8&3&0\\
0&0&-5&-2&0&-3&-4&0&0\\ 0&0&-4&-3&0&-1&-3&-1&0\\ 0&0&4&3&1&3&3&1&0\\
0&0&0&0&0&0&2&-3&0\\ 0&0&0&0&0&0&1&2&0\\ 0&0&0&0&0&0&0&0&3
\end{smallmatrix}\right),
\]
$10^4g^{\mathsf T}=(3329,\,6390,\,4411,\,-611,\,-3579,\,5606,\,4509,\,8550)$;
здесь $\tfrac32A$ --- матрица Грама $E_8$ при $\lambda_1^2=3$, $|\det C|=7203=3\cdot7^4$.
Ячейка имеет $752$ фасеты и $1\,654\,230$ вершин; полнота списка вершин
доказана замыканием по $1$-скелету без предположения простоты, и
\[
\diam^2V_0=\tfrac{1119552292542693}{165294140000000},\qquad
D(\Gamma)^2=7,\qquad d=1{,}0166127\ldots>1 .
\]

В каждом из четырёх случаев \eqref{crit} даёт утверждение теоремы.

\begin{thebibliography}{9}
\bibitem{Soifer} A.~Soifer, \emph{The Mathematical Coloring Book}, Springer,
New York, 2009.
\bibitem{Ivanov2011} Л.~Л.~Иванов, \emph{О хроматических числах $\R^2$ и $\R^3$
с интервалами запрещённых расстояний}, Вестник РУДН, сер. Матем. Информ. Физ.,
2011, №~1, 14--29.
\bibitem{ABPR} A.~Arman, A.~Bondarenko, A.~Prymak, D.~Radchenko, \emph{Upper
bounds on chromatic number of $\mathbb E^n$ in low dimensions}, Electron. J.
Combin., \textbf{31}:2 (2024), P2.35.
\end{thebibliography}

% TODO: место работы обоих авторов (для второго автора --- вместо <<---->>)
\noindent
\parbox[t]{0.5\textwidth}{{\bf Л.\,Л.~Иванов (L.\,L.~Ivanov)}\\
Москва\\
{\it E-mail}: leo.ivanov@gmail.com}%
\parbox[t]{0.5\textwidth}{{\bf Н.\,В.~Глушкова (N.\,V.~Glushkova)}\\
---\\
{\it E-mail}: glushkova.nv@phystech.edu}

\end{document}
